\chapter{Stratified Polyxan--RSVP Correspondence (Discrete $\to$ Continuous)}
\label{chap:28}

The final step in Part~II is the precise translation of discrete Polyxan
structure into continuous RSVP geometry.  
Chapters~25--27 established quines, cohomology, and homotopy.
We now prove that these discrete invariants correspond exactly to
stratified curvature and torsion invariants in the RSVP manifold~$\mathcal{M}$.
This chapter constructs the stratified correspondence functor,
proves its full faithfulness on stable objects,
and identifies its essential image.
It is the completed discrete--to--continuous half of the duality
formalised in Chapter~62 but developed here in the purely mathematical setting.

\section{Stratified Embeddings and Discrete Strata}

Let $\mathcal{G}$ be a finite typed Polyxan hypergraph
and $\mathrm{Inc}(\mathcal{G})$ its incidence complex.
A \emph{stratified embedding}
\[
\iota : \mathrm{Inc}(\mathcal{G}) \hookrightarrow \mathcal{M}
\]
is a smooth, cellular map satisfying:

\begin{enumerate}
\item \textbf{Curvature Matching:}
\[
\mathcal{R}\circ\iota = \mathcal{H}_0[\mathcal{G}],
\]
where $\mathcal{H}_0$ is primitive Polyxan curvature (Ch.~25).

\item \textbf{Tag Entropy Matching:}
\[
S\circ\iota = \mathcal{H}_{\mathrm{tag}}[\mathcal{G}],
\]
where $\mathcal{H}_{\mathrm{tag}}$ is the tag sheaf entropy (Ch.~26).

\item \textbf{Stratum Preservation:}
If a cell $c$ has quine depth $d$,  
then $\iota(c)$ lies entirely in curvature stratum $\mathcal{M}^{(d)}$
(Ch.~13).
\end{enumerate}

Thus every discrete curvature or entropy value is represented
as a continuous geometric invariant.

\section{The Stratified Correspondence Functor}

Define a functor
\[
\mathcal{S} : \mathbf{Polyx}^{\mathrm{typed}}
  \longrightarrow \mathbf{RSVP}^{\mathrm{strat}}
\]
by
\[
\mathcal{S}(\mathcal{G}) := \iota(\mathrm{Inc}(\mathcal{G})),
\]
where $\iota$ is the canonical minimal--energy embedding
(constructed via the embedding flow of Chapter~31).

On morphisms:
a Polyxan rewrite $\mathcal{G}\to\mathcal{H}$
is sent to the unique harmonic extension
of $\iota_{\mathcal{G}}$ into $\iota_{\mathcal{H}}$.

\begin{proposition}[Functoriality]
$\mathcal{S}$ is a functor.
\end{proposition}

\begin{proof}
Composite rewrites correspond to concatenated harmonic extensions.
Curvature and entropy matching are preserved at every stage.
Identity rewrites map to stationary embeddings.
\end{proof}

\section{Full Faithfulness on Quine--Stable Objects}

Let $\mathfrak{Q}$ be any quine--stable hypergraph.

\begin{theorem}[Full Faithfulness]
\label{thm:28-fullfaithful}
For quine--stable objects $\mathfrak{Q}_1,\mathfrak{Q}_2$,
the functor $\mathcal{S}$ induces a bijection
\[
\mathrm{Hom}_{\mathbf{Polyx}}(\mathfrak{Q}_1,\mathfrak{Q}_2)
\;\cong\;
\mathrm{Hom}_{\mathbf{RSVP}}(\mathcal{S}\mathfrak{Q}_1,\mathcal{S}\mathfrak{Q}_2).
\]
\end{theorem}

\begin{proof}
A morphism of quine--stable Polyxan hypergraphs is determined entirely by
its action on self--rewriting atoms.
Under $\mathcal{S}$, each such atom corresponds to a unique harmonic 0--form.
Harmonic forms are uniquely determined by their boundary values.
Thus any morphism in RSVPs between $\mathcal{S}\mathfrak{Q}_1$ and $\mathcal{S}\mathfrak{Q}_2$
corresponds to a unique discrete morphism.
The curvature stratum alignment ensures injectivity.
\end{proof}

\section{Essential Image: Harmonic Stratified Submanifolds}

We now describe the objects in the image of $\mathcal{S}$.

\begin{theorem}[Essential Image]
\label{thm:28-essentialimage}
An RSVP stratified object $X\subseteq\mathcal{M}$
lies in the essential image of $\mathcal{S}$ if and only if:
\begin{enumerate}
\item $X$ has finitely many curvature strata;
\item each stratum is a harmonic submanifold;
\item adjacent strata differ in curvature by exactly $1$;
\item the torsion depth of $X$ is finite.
\end{enumerate}
\end{theorem}

\begin{proof}
(\emph{If})  
Given $X$ with these four properties, construct discrete cells by taking
connected components of each stratum and assigning primitive curvature equal to
stratum depth.
Tag entropy is determined uniquely by harmonicity (Ch.~26).
The resulting hypergraph $\mathcal{G}_X$ satisfies
$\mathcal{S}(\mathcal{G}_X)=X$.

(\emph{Only If})  
If $X=\mathcal{S}(\mathcal{G})$,
then finiteness follows from finiteness of $\mathcal{G}$.
Curvature jumps are exactly primitive curvature increments.
Harmonicity and torsion depth correspond to quine stability
and homotopy height (Ch.~27).
\end{proof}

Thus the stratified correspondence identifies discrete semantic quine structure
with continuous geometric harmonic structure.

\section{Discrete Cohomology = Continuous Harmonic Cohomology}

\begin{theorem}[Cohomology Correspondence]
\label{thm:28-cohom}
For any finite typed $\mathcal{G}$,
\[
H^k_{\mathrm{Polyx}}(\mathcal{G})
  \;\cong\;
\mathcal{H}^k_{\mathrm{harm}}(\mathcal{S}(\mathcal{G})),
\]
where $H^k_{\mathrm{Polyx}}$ is the Čech--de~Rham cohomology (Ch.~26)
and $\mathcal{H}^k_{\mathrm{harm}}$ is the RSVP harmonic cohomology
(Ch.~15).
\end{theorem}

\begin{proof}
Stratum preservation ensures that local Polyxan forms map to local RSVP forms
with identical coboundary operators.
Minimal--energy embeddings guarantee harmonicity of the resulting forms.
\end{proof}

\begin{corollary}[Media Quines Are Harmonic Minimizers]
If $\mathcal{G}$ is a media quine then
$\mathcal{S}(\mathcal{G})$ is a minimal harmonic stratified submanifold.
\end{corollary}

\section{Discrete Homotopy = Continuous Torsion Depth}

Finally, we state the continuous version of the homotopy correspondence
proved discretely in Chapter~27.

\begin{theorem}[Homotopy--Torsion Correspondence]
\label{thm:28-homotopytorsion}
Let $\mathcal{G}$ be a finite Polyxan hypergraph and
$\iota$ the canonical embedding.
Then
\[
h(\mathcal{G})
  \;=\;
\tau\bigl(\iota(\mathcal{G})\bigr),
\]
where $h$ is Polyxan homotopy height (Ch.~27)  
and $\tau$ is RSVP torsion depth (Ch.~14).
\end{theorem}

\begin{proof}
Rewrite spheres correspond to torsion spheres under $\iota$.
Curvature strata track quine depth exactly.
Thus the minimal number of rewrite layers needed for quine stability
equals the minimal torsion layer depth needed for harmonic flattening.
\end{proof}

\bigskip

\begin{center}
\textbf{This completes Chapter~28.}\\[4pt]
The discrete invariants of Polyxan hypergraphs are now fully identified with  
the continuous invariants of RSVP geometry.  
Part~III will use this correspondence to unify the coupled dynamics.
\end{center}

